% Осень26!A38; нумерация внутри столбца без пустых строк.
% Источник: лекция 6, §4.6; лекция 7, «Долги и упражнения»
\ticket{35}{Теорема Леви}
\label{ticket:35}

\materialnote{Есть формулировка и применение; доказательства в лекциях 1--7 нет.}

\begin{theorem}

\emph{монотонная сходимость.} Пусть \((X,\mathcal A,\mu)\) --- пространство с мерой, \(f_n:X\to[0,+\infty]\) --- измеримые функции и \[0\le f_n\le f_{n+1},\qquad f_n(x)\longrightarrow f(x)\quad(x\in X).\] Тогда \[\lim_{n\to\infty}\int_X f_n\,d\mu=\int_X f\,d\mu.\]

\end{theorem}

\textbf{Статус в источниках.} В лекции 6 этот переход использован при вычислении интеграла по взвешенному образу меры. В конце лекции 7 теорема по-прежнему записана в долгах; доказательства в предоставленных материалах нет.
